\documentclass[10pt,a4paper]{amsart}
\usepackage[margin=19mm]{geometry}
\usepackage{amsmath,amssymb,mathtools}
\usepackage[scr=rsfs,cal=euler]{mathalfa}
\usepackage{xfrac}
\usepackage[unicode,hidelinks,bookmarks=false]{hyperref}
\newcommand{\ie}{i.e.\ }
\newcommand{\cf}{cf.\ }
\newcommand{\resp}{respectively\ }
\newcommand{\Subsection}{\subsection}
\providecommand{\nobreakdash}{\nobreak\hbox{-}}
\setlength{\emergencystretch}{2em}
\multlinegap0pt
\usepackage{bm}
\usepackage{bbm}
\usepackage{dsfont}
\usepackage{xspace}
\usepackage{cancel}
\usepackage{mathtools}
\usepackage{amsthm}
\makeatletter
\@ifundefined{equation}{\newtheorem{equation}{Equation}}{}
\@ifundefined{prop}{\newtheorem{prop}[equation]{Proposition}}{}
\@ifundefined{theorem}{\newtheorem{theorem}[equation]{Theorem}}{}
\makeatother
\DeclareMathOperator{\Dom}{Dom}
\DeclareMathOperator{\Exp}{Exp}
\newcommand{\RR}{\mathcal R}
\newcommand{\Z}{\mathbb Z}
\title[Non-degenerate near-parabolic renormalization]{Non-degenerate near-parabolic renormalization}
\author{Alex Kapiamba}
\date{Source: arXiv:2210.06647v4 (CC BY 4.0)}
\begin{document}
\maketitle

\noindent\emph{Source and licence.} Non-degenerate near-parabolic renormalization. Version arXiv:2210.06647v4 (CC BY 4.0), distributed under the Creative Commons Attribution 4.0 International licence, as recorded in the arXiv licence metadata for this exact version and shown on the abstract page https://arxiv.org/abs/2210.06647v4. Full rights notice: licenses/arxiv-2210.06647-CC-BY-4.0.txt.\par\smallskip

\noindent\textit{Editorial scope.} Counted unit: the Theorem whose source label is \texttt{thm:comparison} in the article (source0.tex lines 1167--1170, source label \texttt{thm:comparison}), delivered together with the article's own complete original proof of it (source0.tex lines 1173--1197, one \texttt{proof} environment of 25 lines). No mathematics is written, repaired or re-derived: statement and proof are the article's own text line for line. Required context retained uncounted: lifing the renormalized dynamics (lines 821--834); prop:comparing petals (lines 1143--1145). Every definition, display and label of those lines is retained unchanged. Omitted from this package and not redistributed: 1--820, 835--1142, 1146--1166, 1171--1172, 1198--1979. Nothing inside a retained excerpt is omitted or altered: every retained body is delivered exactly as the article's lines are written, so no omission receipt of any kind accompanies this package. Citations: the retained excerpts carry no citation command at all, so no bibliography entry of the article is transcribed and no reference is redistributed. Reference adaptation: none was needed and none was made; every reference inside the delivered unit resolves inside it, so no unresolved reference or citation is produced. Printed slips kept literal: no printed word, symbol, hypothesis or number of the counted statement or of its proof was altered. Layout and class: the article's own class is \texttt{amsart} with packages including amsmath, amssymb, amsthm, graphicx, hyperref, mathtools, verbatim, geometry, tikz-cd, enumitem, todonotes, biblatex; the wrapper is the lane's pinned \texttt{amsart} editorial wrapper at 10pt on A4, which restates the theorem-like environments the retained text needs and adds the article's own macro definitions that the retained text needs (\texttt{\textbackslash Dom}, \texttt{\textbackslash Exp}, \texttt{\textbackslash RR}, \texttt{\textbackslash Z}), with extra wrapper packages (bm, bbm, dsfont, xspace, cancel, mathtools). The delivered numbering is the unit's own. Assets: no retained span contains a figure environment, a graphics-inclusion command, a PDF-page inclusion, a table or a TikZ picture; no image file is copied into this package, the package declares no asset dependency and no image file is redistributed with it. Licence: version arXiv:2210.06647v4 of this article is distributed under the Creative Commons Attribution 4.0 International licence, as recorded in the arXiv licence metadata for this exact version and shown on the abstract page https://arxiv.org/abs/2210.06647v4. A full rights notice is delivered at licenses/arxiv-2210.06647-CC-BY-4.0.txt. Characters: every retained excerpt is pure ASCII, so the delivered engine, XeLaTeX with its default Latin Modern text font, reports no missing character.
	\begin{prop}\label{lifing the renormalized dynamics}
		Let $z, z'$ be two points in $P_0$ and set $w= \phi_0(w)$, $w'= \phi_0(z')$. If 
		$$H_\pm\circ T_{n- 1/\alpha}(w) = w'$$
		for some integer $n$, then there are integers $m\geq 0$ and $0\leq k < q$ such that either
		\begin{align*}
			g^{(n+m)q+k}(z)&= g^{mq}(z') &&\hspace{-3cm}\text{ if }n\geq 0, \text{ or}\\
			g^{mq+k}(z) & = g^{(m-n)q}(z') && \hspace{-3cm}\text{ if }n\leq 0.
		\end{align*}
		If $z'$ is exiting $P_0$, then $m = 0$. If $|\emph{Im}\,w|$ is sufficiently large, then $m = 0$ and 
		$$k = \begin{cases}
			q_\pm' & \text{ if }\pm\emph{Im}\, w>0,\\
			0 &\text{ if }\mp\emph{Im}\, w< 0.
		\end{cases}$$
	\end{prop}

	\begin{prop}\label{prop:comparing petals}
		We can choose the flowers  so that $\Exp_{s_0}\circ \phi_0^{h_0, f_0}$ maps $P_0^{h_0, g_0}$ to $P_0^{h_1, f_1}$, conjugates $h_0^q$ to $h_1^{q_1}$ there, and satisfies $$\phi_0^{h_0, g_0} = \phi_1^{h_1, f_1}\circ \Exp_{s_0}\circ \phi_0^{h_0, f_0}.$$
	\end{prop}

	\begin{theorem}\label{thm:comparison}
		For any compact $X\in \Dom(f_2)$, if $g_0$ is close to $f_0$ and $h_0$ is close  to $g_0$, then we can choose the  near-parabolic renormalizations so that 
		$h_2= \RR_{g_0}^sh_0$ on $X$.
	\end{theorem}

	\begin{proof}
		Let $H^{f_0}_{s_0}$ and $H^{f_1}_{s_1}$ be horn maps for $f_0$ and $f_1$ respectively. 
		Let $H^{h_0, g_0}$, $H^{h_0, f_0}$,  and $H^{h_1, f_1}$ be horn maps for $h_0$ relative to $g_0$, $h_0$ relative to $f_0$, and $h_1$ relative to $f_1$ respectively. As $\Exp_{s_1}= \Exp_{s}$ semi-conjugates $H^{h_1, f_1}\circ T_{-1/\alpha}$ and $H^{h_0, g_0}\circ T_{-1/\alpha}$ to $h_2$ and $\RR_{g_0}h_0$ respectively, the theorem is a consequence of the following  claim: for any compact $X\subset \Dom(H^{f_1}_{s_1})$, we can choose the horn maps so that 
		$H^{h_1, f_1}-H^{h_0, g_0}$ maps $X$ into $\Z$.
		
		Let $X_1\subset P_{s_1}^{f_1}$ be a compact set such that $f_1^{N_1}(X_1)\subset P_0^{f_1}$ 
		for some integer $N_1\geq 0$. 
		Thus $X_1\subset P_{s_1}^{h_1, f_1}$ and $h_1^{N_1}(X_1)$ is entering $P_0^{h_1, f_1}$ when $h_1$ is close to $f_1$. 
		Assuming the flowers are chosen as in Proposition \ref{prop:comparing petals}, let $X_0$ be the component of $(\Exp_{s_0}\circ \phi_0^{h_0, f_0})^{-1}(X_0)$ contained in $P_0^{h_0, g_0}$.
		Note that it follows from the proof of Proposition \ref{prop:comparing petals} that we can choose the petals so that $P_0^{h_0, g_0}$ is exiting $P_0^{h_0, f_0}$.
		
		
		
		We set $Y_0= \phi_0^{h_0, f_0}(X_0)$, so $\Exp_{s_0}(Y_0)= X_1$. 
		Setting $n_0 = n$, and 
		choosing some integers $j_m$ for all $0\leq m < m_1$, we inductively define $$Y_{m+1} = H^{h_0, f_0}\circ T_{n_0+j_m-1/\alpha_0} (Y_m).$$
		When $h_1$ is sufficiently close to $f_1$ and $h_0$ is sufficiently close to $f_0$, we can choose the integers $j_m$, depending only on $X_1$, so that each $(\phi_0^{h_0, f_0})^{-1}(Y_m)$ is defined and exiting $P_0^{h_0, f_0}$. 
		Moreover, as $h_1^{N_1}(X_1)$ is entering $P_0^{h_1, f_1}$, we can choose the integers so that $(\phi_0^{h_0, f_0})^{-1}(Y_{N_1})$ is entering $P_0^{h_0, g_0}$.
		It then follows from Proposition \ref{lifing the renormalized dynamics} that when $\alpha$ is sufficiently small there is some $$N_0\leq \sum_{m=0}^{m_1-1}(n_0+j_m+1)q_0< \text{Re}(1/3\alpha)q$$
		such that $h_0^{N_0}(X_0)$ is entering $P_0^{h_0, g_0}$
		and $\Exp_{s_0}\circ \phi_0^{h_0, f_0}\circ h_0^{N_0} = h_1^{N_1}$ on $X_0$. 
		Hence for any $z_0\in X_0$, $z_1 = \Exp_s\circ \phi_0^{h_0, f_0}(z_0)$, and $w= \phi_0^{h_0, g_0}(z_0) = \phi_0^{h_1, f_1}(z_1)$, we have that 
		$$H^{h_0, g_0}(w) \equiv \phi_0^{h_0, g_0}\circ h_0^{N_0}(z_0) = \phi_0^{h_1, f_1}\circ h_1^{N_1}(z_1)\equiv H^{h_1, f_1}(w)$$
		modulo $\Z$. For any compact set $X\subset \Dom(H^{f_1}_{s_1})$, we can choose $X_1$ so that $\phi_0^{h_1, f_1}(X_1)+m$ contains $X$ for some integer $m$, which completes the proof.
	\end{proof}

\end{document}
